\chapter{المخرج السريع الوحيد}
% full version: chapters/13_muscles.tex

كل ما تفعله في العالم بسرعة، من كلمة أو نظرة أو خطوة، هو انقباض للعضلات. وللآلة مخرج ثانٍ بطيء، كيميائي، وعنه الفصل السادس عشر. أما هنا فالعضلات.

آخر حلقة في السلسلة عصبونات \nt{ACET}. كل ليف عضلي يصغي إلى عصبون واحد بالضبط من هذا النوع، والعصبون الواحد يتحكم في مجموعة من الألياف، من المئات إلى بضعة آلاف. واللافت أن الوصلة بين العصب والعضلة أوثق وصلة في الجسم. داخل الدماغ تضيّع الوصلات أكثر من نصف الإشارات، أما هنا فكل نقرة تقذف من المادة أضعاف ما يلزم. عند المخرج يكلف الخطأ حركة، فالموثوقية مشتراة بهامش واسع.

\section*{قرع الطبل يصير دفعة}

نقرة واحدة تعني ارتعاشة قصيرة واحدة للعضلة. إذا تباعدت النقرات، ارتعشت العضلة بهزات منفصلة. وإذا تقاربت، اندمجت الارتعاشات في جهد متصل، كما يندمج قرع الطبل في دويّ. العضلة تحوّل تدفق النبضات إلى قوة، كأبسط محوّل رقمي تماثلي.

\pencilfig{113}{%
% twitch = alpha function; the sum of twitches for spikes 1 s and 0.16 s apart (arbitrary units of time)
\begin{scope}[x=0.95cm, y=0.36cm]
\foreach \dx/\t in {0/متباعدة,4.6/متقاربة}{
  \begin{scope}[xshift=\dx cm]
  \draw[pen] (0,0) -- (4,0); \node[font=\small\itshape, text=pcGray, anchor=south] at (2,5.4) {\t};
  \end{scope}}
\draw[penlite=pcRed] plot[smooth] coordinates {(0.00,0.00) (0.05,0.00) (0.10,0.00) (0.15,0.00) (0.20,0.00) (0.25,0.45) (0.30,0.73) (0.35,0.90) (0.40,0.98) (0.45,1.00) (0.50,0.98) (0.55,0.94) (0.60,0.88) (0.65,0.81) (0.70,0.74) (0.75,0.66) (0.80,0.59) (0.85,0.52) (0.90,0.46) (0.95,0.41) (1.00,0.35) (1.05,0.31) (1.10,0.27) (1.15,0.23) (1.20,0.20) (1.25,0.62) (1.30,0.88) (1.35,1.02) (1.40,1.08) (1.45,1.09) (1.50,1.06) (1.55,1.00) (1.60,0.93) (1.65,0.86) (1.70,0.78) (1.75,0.70) (1.80,0.62) (1.85,0.55) (1.90,0.48) (1.95,0.42) (2.00,0.37) (2.05,0.32) (2.10,0.28) (2.15,0.24) (2.20,0.21) (2.25,0.62) (2.30,0.88) (2.35,1.03) (2.40,1.09) (2.45,1.09) (2.50,1.06) (2.55,1.01) (2.60,0.94) (2.65,0.86) (2.70,0.78) (2.75,0.70) (2.80,0.62) (2.85,0.55) (2.90,0.48) (2.95,0.42) (3.00,0.37) (3.05,0.32) (3.10,0.28) (3.15,0.24) (3.20,0.21) (3.25,0.62) (3.30,0.88) (3.35,1.03) (3.40,1.09) (3.45,1.09) (3.50,1.06) (3.55,1.01) (3.60,0.94) (3.65,0.86) (3.70,0.78) (3.75,0.70) (3.80,0.62) (3.85,0.55) (3.90,0.48) (3.95,0.42) (4.00,0.37)};
\foreach \s in {0.20,1.20,2.20,3.20}{\draw[pcBlue, line width=0.8pt] (\s,-0.35) -- (\s,-1.2);}
\begin{scope}[xshift=4.6cm]
\draw[penlite=pcRed] plot[smooth] coordinates {(0.00,0.00) (0.05,0.00) (0.10,0.00) (0.15,0.00) (0.20,0.00) (0.25,0.45) (0.30,0.73) (0.35,0.90) (0.40,1.35) (0.45,1.68) (0.50,1.85) (0.55,2.19) (0.60,2.51) (0.65,2.64) (0.70,2.84) (0.75,3.13) (0.80,3.22) (0.85,3.27) (0.90,3.55) (0.95,3.62) (1.00,3.54) (1.05,3.82) (1.10,3.88) (1.15,3.80) (1.20,3.98) (1.25,4.06) (1.30,3.97) (1.35,4.07) (1.40,4.16) (1.45,4.09) (1.50,4.10) (1.55,4.23) (1.60,4.17) (1.65,4.10) (1.70,4.26) (1.75,4.23) (1.80,4.07) (1.85,4.27) (1.90,4.27) (1.95,4.12) (2.00,4.26) (2.05,4.29) (2.10,4.17) (2.15,4.24) (2.20,4.31) (2.25,4.21) (2.30,4.21) (2.35,4.31) (2.40,4.24) (2.45,4.16) (2.50,4.31) (2.55,4.27) (2.60,4.10) (2.65,4.30) (2.70,4.29) (2.75,4.15) (2.80,4.28) (2.85,4.31) (2.90,4.18) (2.95,4.25) (3.00,4.32) (3.05,4.22) (3.10,4.21) (3.15,4.32) (3.20,4.25) (3.25,4.17) (3.30,4.31) (3.35,4.27) (3.40,4.11) (3.45,3.86) (3.50,3.56) (3.55,3.25) (3.60,2.93) (3.65,2.63) (3.70,2.33) (3.75,2.06) (3.80,1.81) (3.85,1.58) (3.90,1.38) (3.95,1.19) (4.00,1.03)};
\foreach \s in {0.20,0.36,0.52,0.68,0.84,1.00,1.16,1.32,1.48,1.64,1.80,1.96,2.12,2.28,2.44,2.60,2.76,2.92,3.08,3.24}{\draw[pcBlue, line width=0.8pt] (\s,-0.35) -- (\s,-1.2);}
\end{scope}
\node[lbl, text=pcRed, anchor=west] at (1.2,1.9) {ارتعاشات};
\node[lbl, text=pcRed, anchor=south] at (6.6,4.55) {جهد متصل};
\node[lbl, text=pcBlue, anchor=north] at (4.3,-1.35) {نقرات العصبون};
\end{scope}
}{النقرات المتباعدة تعطي ارتعاشات منفصلة للعضلة؛ والمتقاربة تندمج في قوة متصلة أكبر بكثير.}

\section*{قوة بفاصلة عائمة}

للدماغ مقبضان للقوة: كم مرة تنقر العصبونات، وكم مجموعة من الألياف تعمل. تُشغَّل المجموعات بالترتيب، من الأضعف إلى الأقوى، ولا أحد يحسب هذا الترتيب: العصبونات الصغيرة ببساطة أسهل إثارة. والنتيجة أن كل مجموعة جديدة تضيف تقريبًا \emph{النسبة} نفسها إلى القوة الموجودة. حين تمسك بيضة، تُجرَّع القوة بحصص ضئيلة، وحين ترفع حقيبة سفر، بحصص كبيرة. سيرى المهندس هنا عددًا بفاصلة عائمة: الرتبة يحددها عدد المجموعات العاملة، والقيمة الدقيقة يحددها تردد النقرات. والوجه الآخر: كلما قويت الحركة، قلت دقتها. ووفق فرضية مؤثرة، لهذا السبب حركاتنا سلسة: الأمر السلس يعطي أقل تشتت.

\section*{تشد ولا تدفع}

العضلة لا تستطيع إلا أن تشد. لذلك يخدم كل مفصل زوج من العضلات تشدان في اتجاهين متعاكسين: مرة أخرى سلكان بدل علامة الزائد والناقص. فرق قوتيهما يدير المفصل، ومجموعهما يجعله أصلب. حين تحمل فنجانًا ممتلئًا، تشد العضلتين معًا: الدوران نفسه، لكن اليد تصمد للصدمة أفضل.

\pencilfig{13}{\pencilart{13}}{العضلة لا تستطيع إلا أن تشد، لذلك هما اثنتان: فرق القوتين يدير المفصل، ومجموعهما يجعله أصلب.}

\section*{تغذية راجعة متأخرة}

في العضلات مستشعرات للتمطط، وأقصر حلقة تُغلق مباشرة في الحبل الشوكي: شُدّت العضلة فانقبضت هي نفسها، وتعطلت العضلة المضادة خلال ذلك. يعطلها عصبون كابح من نوع \nt{GLYC}: وها قد وجدنا عملًا لهذا النوع.

لكن لكل حلقة تأخيرًا: في الحبل الشوكي عشرات الأجزاء من ألف من الثانية، وعبر العين مئات. والتغذية الراجعة المتأخرة تجعل النظام يتأرجح، كسائق يوجّه المقود وهو يرى الطريق متأخرًا. كلما طالت الحلقة، وجب أن تصحح أبطأ. والحد الذي تبدأ عنده الحلقة الشوكية بالتأرجح نحو ثماني ذبذبات في الثانية. وهذا قريب من تردد الرعشة الدقيقة الموجودة عند كل إنسان سليم، وحلقة التمطط أحد مصادرها المحتملة.

فكيف نتحرك بسرعة ودقة؟ وفق فرضية شائعة: نتنبأ. نحمل نموذجًا داخليًا للجسم ونحسب نتيجة الأمر دون انتظار المستشعرات.

\section*{إيقاع بلا ميقاتية}

المشي والتنفس إيقاعيان، لكنهما لا يحتاجان إلى ميقاتية. مجموعتان من العصبونات تكبح كل منهما الأخرى وتتعبان تدريجيًا تعملان بالتناوب من تلقاء نفسيهما: يتعب الفائز، فيغلب المنافس المستريح، وهكذا مرة بعد مرة. والأمر الثابت الآتي من الأعلى، ”امشِ“، يتحول في موضعه إلى إيقاع ”يسار، يمين“.

\fullversion{13}
